\section{Benefits of Proposed System}

SmartExam delivers nine measurable benefits to the institutions and individuals that adopt it.

\textbf{1. Role-based administration.} In the system, administrators will gain a more centralized approach for user and data management, which defines what each user can and cannot do to prevent misuse of information and functions.

\textbf{2. Structured scheduling.} Teachers can schedule a specific timed exam based on parameters they define such as title, time limitation, permitted list of applications or coding examination parameters such as test cases and instructions \cite{eassess1}.

\textbf{3. Hardware restricted system.} The system locks the user account to a specified machine through hardware ID registration. The system will not allow any student to access the application or run exams on unauthorized hardware, such as a peer's computer or personal machine that student has not used previously.

\textbf{4. ManagedExamSession.} Student client in WPF runs in fullscreen mode to disable Alt+Tab, Task Manager and clipboard access. Policy is enforced technically rather than manually by verbal warnings.

\textbf{5. LiveMonitoring Dashboard.} Invigilator view shows active student sessions in a grid. For each student, there is status indicator for focus, list of active processes, screenshot thumbnails of recent frames, and a focus trajectory. When specific thresholds are exceeded automatically triggers alerts.

\textbf{6. Auditableincident Recording.} Any recorded event (e.g. Loss of student focus, denied app execution, or idle state) is instantly saved to the database, with an association between timestamp and session ID. Records persist beyond the session for audit purposes.

\textbf{7. Submit reliably.} Student automatically autosaves answers several times per minute. The final version, submitted either manually when student or remotely by teacher, will be sent to the server where it will be locked forever.

\textbf{8. AI Supported Grading.} Coding submissions are validated via the tests of the pre-defined environment. Textual and theoretical tests are marked by comparing to previously specified answers using semantics based calculation. Reduces teacher work-load significantly, especially for the numerous tests cases \cite{spp1}.

\textbf{9. Analytical reports.} At the end of exam session administrator and teacher create one report from various statistics such as submission list, attendance and flagging list. The reports offer hyperlinking functionality to view respective records of incident/submission \cite{sommerville}.
